% Introduction générale en cinq mouvements (trame, étape 4) : contexte large, problème réel, faille, objectif, contribution.
\chapter{Introduction générale}\label{ch:intro}

% 1. Contexte large
Un modèle de langage ne sait pas tout. Pour répondre à une question factuelle rare ou composée de plusieurs étapes, on lui donne accès à un moteur de recherche \citep{lewis2020retrieval}, et la difficulté se déplace : il faut décider quoi chercher, quand s'arrêter et comment se servir de ce que l'on trouve. Deux façons de faire coexistent. Dans la première, un \emph{harnais} écrit par le programmeur fixe la règle : chercher une fois avant de répondre, alterner pensée et action \citep{yao2022react}, ou chercher après chaque étape de raisonnement \citep{trivedi2022interleaving}. Dans la seconde, le modèle apprend sa propre règle par renforcement, en étant récompensé quand sa réponse finale est juste ; Search-R1 \citep{jin2025search} en est l'exemple le plus étudié.

% 2. Problème réel et enjeu
Le choix entre les deux n'est pas seulement académique. Entraîner un agent par renforcement coûte cher : il faut générer des milliers de trajectoires à chaque étape d'entraînement, avec un moteur de recherche dans la boucle. Des entreprises en font pourtant leur métier et vendent cet entraînement comme un moyen d'obtenir de meilleurs agents. Or il existe une voie bien moins coûteuse : collecter les bonnes trajectoires d'un modèle et l'entraîner à les imiter, par apprentissage supervisé. Si le renforcement ne fait que rendre plus fréquentes des réponses que le modèle trouvait déjà de temps en temps, cette voie pourrait en obtenir l'essentiel. S'il rend possibles des réponses que le modèle ne trouvait jamais, elle ne le peut pas.

% 3. Faille
Cette question, « fiabiliser ou élargir ? », est vive pour les modèles de raisonnement sans outil depuis que \citet{yue2025does} ont montré qu'à grand nombre d'essais le modèle de départ rattrape le modèle entraîné. Elle reste presque ouverte pour les agents. Deux études seulement la posent \citep{zhai2026does,zhu2026what}, sur des agents que leurs auteurs ont entraînés eux-mêmes, et sans le test statistique apparié dont \citet{yang2026crossovers} a montré qu'il fait disparaître la plupart des écarts publiés. Personne n'a mesuré ce que l'entraînement a apporté aux agents de recherche publics, face à leur propre modèle de départ. Notre première étude (chapitre~\ref{ch:partie1}) avait la même faiblesse : elle comparait des agents Search-R1 issus d'un modèle de base à des harnais construits sur la version Instruct du même modèle, si bien que l'écart mêlait règle de recherche et modèle.

% 4. Objectif
Ce mémoire pose donc deux questions. La première, traitée au chapitre~\ref{ch:partie1}, compare à conditions égales des agents entraînés et des harnais fixes, au premier essai. La seconde, cœur du mémoire (chapitre~\ref{ch:partie2}), compare chaque agent Search-R1 entraîné à partir de Qwen2.5-3B-Instruct à ce même modèle : à modèle de départ identique, le renforcement élargit-il l'ensemble des questions que l'agent finit par résoudre, ou augmente-t-il seulement sa réussite au premier essai ? Cet élargissement apparaît-il quand la question oblige à enchaîner les recherches ? Et l'avantage vient-il de la recherche ou de la lecture de ce qui est trouvé ?

% 5. Contribution
Le mémoire apporte trois éléments. D'abord, la première mesure appariée de la frontière d'agents de recherche publics face à leur modèle de départ, avec un plan d'analyse gelé et publié avant les données. Ensuite, une séparation expérimentale de la recherche et de la lecture, par trois mesures indépendantes qui se recoupent. Enfin, une correction de notre propre conclusion initiale : à modèle de départ identique, l'avantage de l'agent tient à sa recherche, et sa meilleure lecture est masquée par ses refus de répondre. Le chapitre~\ref{ch:partie1} présente l'étude initiale, le chapitre~\ref{ch:partie2} l'étude appariée ; la discussion générale (chapitre~\ref{ch:discussion}) les relit ensemble.
